
En
effet, soient \(S'\) l'espace affine \(\mathbb A_k^N\) dont les \(k\)-points correspondent aux suites \(\sigma\), et \(T\) l'ouvert de \(S\) correspondant aux suites \(F\)-régulières. Pour \(i\) fixé, les faisceaux \(F_{\sigma,i}(-1)\) qui correspondent aux \(k\)-points de \(T\) sont contenus dans une famille plate sur \(T\), et par hypothèse \(T\) est non-vide. Donc (i) résulte de la propriété de semi-continuité supérieure de la fonction $\sigma\longmapsto h^0(F_{\sigma,i}(-1))$ (EGA III, 7.7.5). Enfin, comme toute suite ``assez générale'' est \(G\)-régulière et telle que \(G_{\sigma,i}\) soit un sous-faisceau de \(F_{\sigma,i}\) pour tout \(i\), (ii) résulte de (i).

\textbf{Lemme 1.7.} \emph{Soit \(0\to F(1)\to F\to G\to0\) une suite exacte sur \(X\). Si le polynôme de Hilbert de \(F\) est écrit sous la forme}
\[
 \chi(F(n))=\sum_{i=0}^{r}a_i\binom{n+i}{i},
\]
\emph{alors}
\[
 \chi(G(n))=\sum_{i=0}^{r-1}a_{i+1}\binom{n+i}{i}.
\]
En effet, \(\chi(G(n))=\chi(F(n))-\chi(F(n-1))=\sum_{i=0}^{r}a_i\left[\binom{n+i}{i}-\binom{n-1+i}{i}\right]\); donc \(\chi(G(n))=\sum_{i=1}^{r}a_i\binom{n+i-1}{i-1}\), d'où l'assertion.

\textbf{Proposition 1.8.} \emph{Soient \(F\) un \((b)\)-faisceau sur \(X\), \(s=\operatorname{dim{.}supp}(F)\) et \(\chi(F(n))=\sum_{i=0}^{s}a_i\binom{n+i}{i}\) le polynôme de Hilbert de \(F\). Alors}~:
\begin{enumerate}[label=(\roman*)]
\item \emph{Pour \(n\ge -1\), \(h^0(F(n))\le \sum_{i=0}^{s}b_i\binom{n+i}{i}\).}
\item \emph{\(a_s\le b_s\) et \(F\) est aussi un \((b_0,\ldots,b_{s-1},a_s)\)-faisceau.}
\end{enumerate}
En effet, raisonnons par récurrence sur \(s\). Lorsque \(s=0\), c'est clair car \(a_0=h^0(F)=h^0(F(-1))\le b_0\). Lorsque \(s\ge1\), on a une suite exacte
\[
0\longrightarrow F(-1)\longrightarrow F\longrightarrow G\longrightarrow0
\]
